\ifdefined\gkver\else\def\gkver{student}\fi
\documentclass[\gkver]{gaokaozhenti}
\begin{document}

\chapter{2014年天津理综}
%% paperType: 理综物理部分

\begin{enumerate}
\item
%% number: 1
%% paperName: 2014年天津理综第 1 题 6 分
%% typeId: 1
%% type: 单选题
%% chapter: 直线运动
%% point: 运动图象
%% method: 无
%% score: 6
%% degree: 650
%% duplicateId: 0
%% body:
质点作直线运动的速度—时间图像如图所示，该质点\xzanswer{D}
\onepicture[align=right, w=5.0cm]{figs/01.png}
\fourchoices[answer=D]
{在第1秒末速度方向发生了改变}
{在第2秒末加速度方向发生了改变}
{在前2秒内发生的位移为零}
{第3秒末和第5秒的位置相同}
%% answer: D
%% memo:
\memoanswer{
由速度图象可知：速度的正负代表了运动的方向，A 错误；图线的斜率代表了加速度的大小及方向，B 错误；图线与时间轴围成的图形的面积代表了物体的位移，C 错误，D 正确。
}

\item
%% number: 2
%% paperName: 2014年天津理综第 2 题 6 分
%% typeId: 1
%% type: 单选题
%% chapter: 电场
%% point: 带电粒子的平衡
%% method: 无
%% score: 6
%% degree: 650
%% duplicateId: 0
%% body:
如图所示，电路中 $R_{1}$、$R_{2}$ 均为可变电阻，电源内阻不能忽略，平行板电容器 C 的极板水平放置。闭合电键 S，电路达到稳定时，带电油滴悬浮在两板之间静止不动。如果仅改变下列某一个条件，油滴仍能静止不动的是\xzanswer{B}
\onepicture[align=right, w=4.5cm]{figs/02.png}
\fourchoices[answer=B]
{增大 $R_{1}$ 的阻值}
{增大 $R_{2}$ 的阻值}
{增大两板间的距离}
{断开电键 S}
%% answer: B
%% memo:
\memoanswer{
开始时刻油滴静止状态，说明受到的电场力与重力等大反向。

增大 $R_{1}$ 的阻值，电容器两端电压变大，根据 $E=\frac{U}{d}$ 可知场强变大，电场力变大，油滴向上运动，A 项错误；从电路图可以看出，电阻 $R_{2}$ 对电容器两端的电压不能起到调节作用，所以改变其大小，不能影响电容器两端的电压，故 B 项正确；增大两极板间距离，根据 $E=\frac{U}{d}$，可知场强减小，电场力减小，所以油滴向下运动，C 项错误；断开电键 S，电容器会放电，所以电场强度会变小，因此油滴不会静止，D 项错误。
}

\item
%% number: 3
%% paperName: 2014年天津理综第 3 题 6 分
%% typeId: 1
%% type: 单选题
%% chapter: 曲线运动
%% point: 人造地球卫星
%% method: 无
%% score: 6
%% degree: 450
%% duplicateId: 0
%% body:
研究表明，地球自转在逐渐变慢，3亿年前地球自转的周期约为22小时。假设这种趋势会持续下去，地球的其他条件都不变，未来人类发射的地球同步卫星与现在相比\xzanswer{A}
\fourchoices[answer=A]
{距地球的高度变大}
{向心加速度变大}
{线速度变大}
{加速度变大}
%% answer: A
%% memo:
\memoanswer{
同步卫星与地球的自转周期相同，若地球的自转周期变大，那么同步卫星的角速度将变小，D 项错误；根据万有引力提供向心力，有：$G\frac{Mm}{(R+h)^{2}}=m\omega^{2}(R+h)$，角速度变小，同步卫星离地高度将变大，A 项正确；由 $a=G\frac{M}{(R+h)^{2}}$ 可知，同步卫星的向心加速度将减小，B 项错误；由 $v=\sqrt{\frac{GM}{R+h}}$ 可知，线速度变小，C 项错误。
}

\item
%% number: 4
%% paperName: 2014年天津理综第 4 题 6 分
%% typeId: 1
%% type: 单选题
%% chapter: 电场
%% point: 带电粒子的偏转
%% method: 无
%% score: 6
%% degree: 450
%% duplicateId: 0
%% body:
如图所示，平行金属板 A、B 水平正对放置，分别带等量异号电荷。一带电微粒水平射入板间，在重力和电场力共同作用下运动，轨迹如图中虚线所示，那么\xzanswer{C}
\onepicture[align=right, w=4.0cm]{figs/04.png}
\fourchoices[answer=C]
{若微粒带正电荷，则 A 板一定带正电荷}
{微粒从 M 点运动到 N 点电势能一定增加}
{微粒从 M 点运动到 N 点动能一定增加}
{微粒从 M 点运动到 N 点机械能一定增加}
%% answer: C
%% memo:
\memoanswer{
根据微粒的运动轨迹可知其合力竖直向下，重力和电场力的大小关系不能确定，所以电场力的方向无法确定，电场力做功无法判断，所以电势能的变化无法确定，AB 项错误。微粒的合力一定做正功，所以微粒的动能一定增加，C 项正确；电场力做功无法判断，所以机械能的变化无法确定，D 项错误。
}

\item
%% number: 5
%% paperName: 2014年天津理综第 5 题 6 分
%% typeId: 1
%% type: 单选题
%% chapter: 机械波
%% point: 波长频率波速
%% method: 无
%% score: 6
%% degree: 450
%% duplicateId: 0
%% body:
平衡位置处于坐标原点的波源 S 在 $y$ 轴上振动，产生频率 $50\UHz$ 的简谐波向 $x$ 轴正、负两个方向传播，波速均为 $100\Ums$。平衡位置在 $x$ 轴上的 P、Q 两个质点随波源振动着，P、Q 的 $x$ 轴坐标分别为 $x_{P}=3.5\Um$、$x_{Q}=-3\Um$。当 S 位移为负且向 $-y$ 方向运动时，P、Q 两质点的\xzanswer{D}
\fourchoices[answer=D]
{位移方向相同、速度方向相反}
{位移方向相同、速度方向相同}
{位移方向相反、速度方向相反}
{位移方向相反、速度方向相同}
%% answer: D
%% memo:
\memoanswer{
本题是对机械振动和机械波的综合考查，根据波速和频率之间的关系 $v=\lambda f$ 得该波的波长为 $2\Um$。可知，当 S 经过平衡位置向负方向运动时，距离 S 为 1.75 个波长的 P 点从最大的负位移处向平衡位置运动，距离 S 为 1.5 个波长的 Q 点经平衡位置向正方向运动，所以两质点的位移方向相反，速度方向相同，D 正确。
}

\item
%% number: 6
%% paperName: 2014年天津理综第 6 题 6 分
%% typeId: 2
%% type: 多选题
%% chapter: 光学
%% point: 光的电磁说
%% method: 无
%% score: 6
%% degree: 650
%% duplicateId: 0
%% body:
下列说法正确的是\xzanswer{BD}
\fourchoices[answer=BD]
{玻尔对氢原子光谱的研究导致原子的核式结构模型的建立}
{可利用某些物质在紫外线照射下发射出荧光来设计防伪措施}
{天然放射现象中产生的射线都能在电场或磁场中发生偏转}
{观察者与波源互相远离时接收到波的频率与波源频率不同}
%% answer: BD
%% memo:
\memoanswer{
卢瑟福通过 $\alpha$ 散射实验提出了原子的核式结构，A 项错误；紫外线具有荧光作用，可以用来设计防伪措施，B 项正确；天然放射现象中有三种射线，其中 $\gamma$ 射线不带电，在电磁场中均不偏转，C 项错误；根据多普勒效应，当观察者与波源互相远离时，观察者接收到的频率小于波源发出的频率，D 项正确。
}

\item
%% number: 7
%% paperName: 2014年天津理综第 7 题 6 分
%% typeId: 2
%% type: 多选题
%% chapter: 交流电
%% point: 交流电
%% method: 无
%% score: 6
%% degree: 650
%% duplicateId: 0
%% body:
如图 1 所示，在匀强磁场中，一矩形金属线圈两次分别以不同的转速，绕与磁感线垂直的轴匀速转动，产生的交变电动势图象如图 2 中曲线 $a$、$b$ 所示，则\xzanswer{AC}
\onepicture[align=right, w=13.0cm]{\includesvg{figs/07}}
\fourchoices[answer=AC]
{两次 $t=0$ 时刻线圈平面与中性面重合}
{曲线 $a$、$b$ 对应的线圈转速之比为 $2:3$}
{曲线 $a$ 表示的交变电动势频率为 $25\UHz$}
{曲线 $b$ 表示的交变电动势有效值为 $10\UV$}
%% answer: AC
%% memo:
\memoanswer{
从图像可看出，从金属线圈旋转至中性面时开始计时，曲线 $a$ 表示的交变电动势周期为 $4\times10^{-2}\Us$，曲线 $b$ 表示的交变电动势的周期为 $6\times10^{-2}\Us$，所以 A、C 正确，B 错误；由 $E_{\mathrm{m}}=NBS\omega$ 可知，$\frac{E_{\mathrm{ma}}}{E_{\mathrm{mb}}}=\frac{\omega_{a}}{\omega_{b}}=\frac{T_{b}}{T_{a}}=\frac{3}{2}$，故 $E_{\mathrm{mb}}=\frac{2}{3}E_{\mathrm{ma}}=10\UV$，曲线 $b$ 表示的交流电动势的有效值为 $5\sqrt{2}\UV$，D 错误。
}

\item
%% number: 8
%% paperName: 2014年天津理综第 8 题 6 分
%% typeId: 2
%% type: 多选题
%% chapter: 光学
%% point: 光的折射
%% method: 无
%% score: 6
%% degree: 450
%% duplicateId: 0
%% body:
一束由两种频率不同的单色光组成的复色光从空气射入玻璃三棱镜后，出射光分成 a、b 两束，如图所示，则 a、b 两束光\xzanswer{AB}
\onepicture[align=right, w=4.2cm]{figs/08.png}
\fourchoices[answer=AB]
{垂直穿过同一平板玻璃，a 光所用的时间比 b 光长}
{从同种介质射入真空发生全反射时，a 光的临界角比 b 光小}
{分别通过同一双缝干涉装置，b 光形成的相邻亮条纹间距小}
{若照射同一金属装置都能发生光电效应，b 光照射时逸出的光电子最大初动能大}
%% answer: AB
%% memo:
\memoanswer{
根据光路图可知 a 光的偏转角较大，说明 a 光的折射率较大，根据 $v=\frac{c}{n}$，可知在同一介质中，a 光所用的时间较长，A 项正确；由 $\sin C=\frac{1}{n}$ 可知，折射率较大的 a 光，其临界角较小，B 项正确；根据干涉条纹间距离：$\Delta x=\frac{L}{d}\lambda$，由折射率的大小关系可知，a 光的波长较小，所以用同一装置做实验时，条纹间距离较小，C 项错误；再由光电效应方程，有：$E_{\mathrm{km}}=h\nu-W_{\text{逸}}$，同一金属的逸出功相同，频率较大的 a 光照射出的光电子的最大初动能较大，D 项错误。
}

\item
%% number: 9
%% paperName: 2014年天津理综第 9 题 4 分
%% typeId: 4
%% type: 实验题
%% chapter: 电路
%% point: 电路实验
%% method: 无
%% score: 4
%% degree: 450
%% duplicateId: 0
%% body:
\textbf{（1）}半径为 $R$ 的水平圆盘绕过圆心 O 的竖直轴转动，A 为圆盘边缘上一点，在 O 的正上方有一个可视为质点的小球以初速度 $v$ 水平抛出时，半径 OA 方向恰好与 $v$ 的方向相同，如图所示。若小球与圆盘只碰一次，且落在 A 点，重力加速度为 $g$，则小球抛出时距 O 的高度 $h=$\tkanswer[2.4cm]{$\frac{gR^{2}}{2v^{2}}$}，圆盘转动的角速度大小 $\omega=$\tkanswer[3.2cm]{$\frac{2n\pi v}{R}$（$n\in\mathrm{N}^{+}$）}。

\onepicture[align=right, w=4.4cm]{figs/09a.png}

\textbf{（2）}某同学把附有滑轮的长木板平放在实验桌面上，将细绳一端拴在小车上，另一端绕过定滑轮，挂上适当的钩码使小车在钩码的牵引下运动，以此定量研究绳拉力做功与小车动能变化的关系。此外还准备了打点计时器及配套的电源、导线、复写纸、纸带、小木块等。组装的实验装置如图所示。

\onepicture[align=right, w=8.4cm]{figs/09b.jpg}

\begin{steps}
	\item
	若要完成该实验，必须的实验器材还有哪些\tkanswer[4.4cm]{刻度尺、天平（包括砝码）}。

	\item
	实验开始时，他先调节木板上定滑轮的高度，使牵引小车的细绳与木板平行。他这样做的目的是下列的哪个\tkanswer[1.2cm]{D}（填字母代号）

	（A）避免小车在运动过程中发生抖动

	（B）可使打点计时器在纸带上打出的点迹清晰

	（C）可以保证小车最终能够实现匀速直线运动

	（D）可在平衡摩擦力后使细绳拉力等于小车受的合力

	\item
	平衡摩擦后，当他用多个钩码牵引小车时，发现小车运动过快，致使打出的纸带上点数太少，难以选到合适的点计算小车的速度。在保证所挂钩码数目不变的条件下，请你利用本实验的器材提出一个解决方法：\tkanswer[4.4cm]{可在小车上加适量的砝码}。

	\item
	他将钩码重力做的功当做细绳拉力做的功，经多次实验发现拉力做功总是要比小车动能增量大一些，这一情况可能是下列哪些原因造成的\tkanswer[1.2cm]{CD}（填字母代号）。

	（A）在接通电源的同时释放了小车

	（B）小车释放时离打点计时器太近

	（C）阻力未完全被小车重力沿木板方向的分力平衡掉

	（D）钩码做匀加速运动，钩码重力大于细绳拉力
\end{steps}

\textbf{（3）}现要测量一个未知电阻 $R_{x}$ 的阻值，除 $R_{x}$ 外可用的器材有：

多用的电表（仅可使用欧姆挡）；

一个电池组 $E$（电动势 $6\UV$）；

一个滑动变阻器 $R$（$0\sim20\UO$，额定电流 $1\UA$）；

两个相同的电流表 G（内阻 $R_{g}=1000\UO$，满偏电流 $I_{g}=100\UuA$）；

两个标准电阻（$R_{1}=29000\UO$，$R_{2}=0.1\UO$）；

一个电键 S、导线若干。

①为了设计电路，先用多用电表的欧姆挡粗测未知电阻，采用“$\times10$”挡，调零后测量该电阻，发现指针偏角非常大，最后几乎紧挨满偏刻度停下来，下列判断和做法正确的是\tkanswer[1.2cm]{AC}（填字母代号）。

（A）这个电阻阻值很小，估计只有几欧姆

（B）这个电阻阻值很大，估计有几千欧姆

（C）如需进一步测量可换“$\times1$”挡，调零后测量

（D）如需进一步测量可换“$\times1\mathrm{k}$”挡，调零后测量

②根据粗测的判断，设计一个测量电路，要求测量尽量准确并使电路能耗较小，画出实验电路图，并将各元件字母代码标在该元件的符号旁。
%% answer: 
%% memo:
\memoanswer{
（1）小球抛出后，水平方向做匀速直线运动，又因为只与圆盘碰撞一次，有：

$R=vt$，$h=\frac{1}{2}gt^{2}$，得 $h=\frac{gR^{2}}{2v^{2}}$；根据圆周运动的周期性，可知两者相撞时，圆盘转动的圈数为整数，有：$\frac{R}{v}=n\frac{2\pi}{\omega}$（$n\in\mathrm{N}^{+}$），得 $\omega=\frac{2n\pi v}{R}$（$n\in\mathrm{N}^{+}$）。

（2）①实验需要测量小车的位移，所以需要刻度尺；另外需要表示小车动能的变化，所以需要小车的质量，因此需要带有砝码的天平；②因为实验要测量拉力对小车所做功与动能的关系，所以实验前需要平衡摩擦力，D 项正确；③纸带上点数过少，说明小车的加速度较大，应适当减小加速度，因此需要在小车上加适量的砝码来达到增大质量，减小加速度的目的；④拉力做功大于小车动能的增量，原因可能是由于小车运动过程中受到的阻力没有被完全平衡；另一个原因是钩码的重力做功，没有完全转化为小车的动能，钩码自身的动能也在增大，CD 项正确。

（3）①欧姆表的指针偏转非常大，说明被测电阻的阻值很小，需要换较小的挡位，因为此时已经选用了 $\times10$ 的挡，所以应该换用 $\times1$ 的欧姆挡，AC 项正确；②此实验器材中提供的电流表都不能直接使用，所以需要对其进行改装，因此需要串联一个较大的电阻 $R_{1}$ 改装成电压表，并联一个较小的电阻 $R_{2}$ 改装成较大量程的电流表；又因为被测电阻较小，所以电流表的分压作用显著，因此需要采用电流表的外接，实验电路如图。

\onepicture[w=7.4cm]{figs/09_an.jpg}
}

\item
%% number: 10
%% paperName: 2014年天津理综第 10 题 16 分
%% typeId: 6
%% type: 计算题
%% chapter: 运动和力
%% point: 动量守恒定律
%% method: 无
%% score: 16
%% degree: 650
%% duplicateId: 0
%% body:
如图所示，水平地面上静止放置一辆小车 A，质量 $m_{A}=4\Ukg$，上表面光滑，小车与地面间的摩擦力极小，可以忽略不计。可视为质点的物块 B 置于 A 的最右端，B 的质量 $m_{B}=2\Ukg$。现对 A 施加一个水平向右的恒力 $F=10\UN$，A 运动一段时间后，小车左端固定的挡板与 B 发生碰撞，碰撞时间极短，碰后 A、B 粘合在一起，共同在 $F$ 的作用下继续运动，碰撞后经时间 $t=0.6\Us$，二者的速度达到 $v_{t}=2\Ums$。求：
\begin{enumerate}
	\item
	A 开始运动时加速度 $a$ 的大小；
	\item
	A、B 碰撞后瞬间的共同速度 $v$ 的大小；
	\item
	A 的上表面长度 $l$。
\end{enumerate}
\onepicture[align=right, w=8.4cm]{figs/10.png}
%% answer: 
\jdanswer{
\begin{enumerate}
	\item
	$2.5\Umsq$

	\item
	$1\Ums$

	\item
	$0.45\Um$

\end{enumerate}
}
%% memo:
\memoanswer{
（1）以 A 为研究对象，根据牛顿第二定律有

$F=m_{A}a$

代入数据解得

$a=2.5\Umsq$

（2）对 A、B 碰后共同运动 $t=0.6\Us$ 的过程，由动量定理得

$Ft=(m_{A}+m_{B})v_{1}-(m_{A}+m_{B})v$

代入数据解得

$v=1\Ums$

（3）设碰前 A 的速度为 $v_{A}$，由动量守恒

$m_{A}v_{A}=(m_{A}+m_{B})v$

A 从开始运动到与 B 碰撞前

$Fl=\frac{1}{2}m_{A}v_{A}^{2}$

代入数据解得

$l=0.45\Um$
}

\item
%% number: 11
%% paperName: 2014年天津理综第 11 题 18 分
%% typeId: 6
%% type: 计算题
%% chapter: 电磁感应
%% point: 电磁感应能量分析
%% method: 无
%% score: 18
%% degree: 450
%% duplicateId: 0
%% body:
如图所示，两根足够长的平行金属导轨固定在倾角 $\theta=30^{\circ}$ 的斜面上，导轨电阻不计，间距 $L=0.4\Um$。导轨所在空间被分成区域Ⅰ和Ⅱ，两区域的边界与斜面的交线为 MN，Ⅰ中的匀强磁场方向垂直斜面向下，Ⅱ中的匀强磁场方向垂直斜面向上，两磁场的磁感应强度大小均为 $B=0.5\UT$。在区域Ⅰ中，将质量 $m_{1}=0.1\Ukg$、电阻 $R_{1}=0.1\UO$ 的金属条 ab 放在导轨上，ab 刚好不下滑。然后，在区域Ⅱ中将质量 $m_{2}=0.4\Ukg$、电阻 $R_{2}=0.1\UO$ 的光滑导体棒 cd 置于导轨上，由静止开始下滑，cd 在滑动过程中始终处于区域Ⅱ的磁场中，ab、cd 始终与导轨垂直且两端与导轨保持良好接触，取 $g=10\Umsq$，问：
\begin{enumerate}
	\item
	cd 下滑的过程中，ab 中的电流方向；
	\item
	ab 将要向上滑动时，cd 的速度 $v$ 多大；
	\item
	从 cd 开始下滑到 ab 刚要向上滑动的过程中，cd 滑动的距离 $x=3.8\Um$，此过程中 ab 上产生的热量 $Q$ 是多少。
\end{enumerate}
\onepicture[align=right, w=6.2cm]{figs/11.png}
%% answer: 
\jdanswer{
\begin{enumerate}
	\item
	由 a 流向 b

	\item
	$5\Ums$

	\item
	$1.3\UJ$

\end{enumerate}
}
%% memo:
\memoanswer{
（1）由 a 流向 b

（2）开始放置 ab 刚好不下滑时，ab 所受摩擦力为最大静摩擦力，设其为 $F_{\max}$

有

$F_{\max}=m_{1}g\sin\theta$

设 ab 刚好要上滑时，cd 棒的感应电动势为 $E$，由法拉第电磁感应定律有

$E=BLv$

设电路中的感应电流为 $I$，由闭合电路的欧姆定律

$I=\frac{E}{R_{1}+R_{2}}$

设 ab 所受安培力为 $F_{\text{安}}$，有

$F_{\text{安}}=ILB$

此时 ab 受到的最大静摩擦力方向沿斜面向下，由平衡条件有

$F_{\text{安}}=m_{1}g\sin\theta+F_{\max}$

代入数据解得

$v=5\Ums$

（3）设 cd 棒的运动过程中电路中产生的总热量为 $Q_{\text{总}}$，由能量守恒有

$m_{2}g\sin\theta\cdot x=Q_{\text{总}}+\frac{1}{2}m_{2}v^{2}$

又

$Q=\frac{R_{1}}{R_{1}+R_{2}}Q_{\text{总}}$

解得

$Q=1.3\UJ$
}

\item
%% number: 12
%% paperName: 2014年天津理综第 12 题 20 分
%% typeId: 6
%% type: 计算题
%% chapter: 磁场
%% point: 带电粒子在磁场中的运动
%% method: 无
%% score: 20
%% degree: 450
%% duplicateId: 0
%% body:
同步加速器在粒子物理研究中有重要的应用，其基本原理简化为如图所示的模型。M、N 为两块中心开有小孔的平行金属板。质量为 $m$、电荷量为 $+q$ 的粒子 A（不计重力）从 M 板小孔飘入板间，初速度可视为零。每当 A 进入板间，两板的电势差变为 $U$，粒子得到加速，当 A 离开 N 板时，两板的电荷量均立即变为零。两板外部存在垂直纸面向里的匀强磁场，A 在磁场作用下做半径为 $R$ 的圆周运动，$R$ 远大于板间距离。A 经电场多次加速，动能不断增大，为使 $R$ 保持不变，磁场必须相应地变化。不计粒子加速时间及其做圆周运动产生的电磁辐射，不考虑磁场变化对粒子速度的影响及相对论效应。求：
\begin{enumerate}
	\item
	A 运动第 1 周时磁场的磁感应强度 $B_{1}$ 的大小。
	\item
	在 A 运动第 $n$ 周的时间内电场力做功的平均功率 $P_{n}$。
	\item
	若有一个质量也为 $m$、电荷量为 $+kq$（$k$ 为大于 1 的整数）的粒子 B（不计重力）与 A 同时从 M 板小孔飘入板间，A、B 初速度均为零，不计两者间的相互作用，除此之外，其他条件均不变。下图中虚线、实线分别表示 A、B 的运动轨迹。在 B 的轨迹半径远大于板间距离的前提下，请指出哪个图能定性地反映 A、B 的运动轨迹，并经推导说明理由。
\end{enumerate}
\twopicture[labelA=2014天津理综12a, labelB=2014天津理综12b, align=right, wA=5.0cm, wB=12.5cm]
{figs/12a.png}
{figs/12b.png}
%% answer: 
\jdanswer{
\begin{enumerate}
	\item
	$B_{1}=\frac{1}{R}\sqrt{\frac{2mU}{q}}$

	\item
	$\overline{P}=\frac{qU}{\pi R}\sqrt{\frac{nqU}{2m}}$

	\item
	图 A。

\end{enumerate}
}
%% memo:
\memoanswer{
（1）设 A 经电场第 1 次加速后速度为 $v_{1}$，由动能定理得：$qU=\frac{1}{2}mv_{1}^{2}-0$

A 在磁场中做匀速圆周运动，所受洛伦兹力充当向心力 $qv_{1}B_{1}=\frac{mv_{1}^{2}}{R}$

由上两式得：$B_{1}=\frac{1}{R}\sqrt{\frac{2mU}{q}}$

（2）设 A 经 n 次加速后的速度为 $v_{n}$，由动能定理，有 $nqU=\frac{1}{2}mv_{n}^{2}-0$

设 A 做第 n 次圆周运动的周期为 $T_{n}$，有 $T_{n}=\frac{2\pi R}{v_{n}}$

设在 A 运动第 n 周的时间内电场力做功为 $W_{n}$，则 $W_{n}=qU$

在该段时间内电场力做功的平均功率为：$\overline{P}=\frac{W_{n}}{T_{n}}$

综上解得：$\overline{P}=\frac{qU}{\pi R}\sqrt{\frac{nqU}{2m}}$

（3）A 图能定性地反映 A、B 运动的轨迹。

A 经过 n 次加速后，设其对应的磁感应强度为 $B$。联立上式并分别应用 A、B 的数据得：

$T_{n}=\frac{2\pi m}{qB_{n}}$

$T'=\frac{2\pi m}{kqB_{n}}=\frac{T_{n}}{k}$

由上知，$T_{n}$ 是 $T'_{n}$ 的 k 倍，所以 A 每绕行 1 周，B 就绕行 k 周，由于电场只在 A 通过时存在，故 B 仅在与 A 同时进入电场时才被加速。

经 n 次加速后，A、B 的速度分别为：$v_{n}$、$v'_{n}$，考虑上式

$v_{n}=\sqrt{\frac{2nqU}{m}}$，$v'_{n}=\sqrt{\frac{2nkqU}{m}}=\sqrt{k}v_{n}$

由题设条件并考虑到上式，对 A 有：$Tv_{n}=2\pi R$

设 B 的轨迹半径为 $R'$，有：$T'v'_{n}=2\pi R'$

比较上述两式得：$R'=\frac{R}{\sqrt{k}}$

上式表明，运动过程中 B 的轨迹半径始终不变。

由以上分析可知，两粒子运动的轨迹如图 A 所示。
}

\end{enumerate}

\end{document}
